% =====================================================================
% Preamble per "Basi di Dati - Modulo 2"
% Adattato dal template Strutture Discrete
% =====================================================================

\usepackage{name}

\raggedbottom
\setlength{\headheight}{14.5pt}

\usepackage{pgfplots}
\pgfplotsset{compat=1.18}
\usepackage{amsthm}
\usepackage{mdframed}
\usepackage[most]{tcolorbox}
\usepackage{listings}
\usepackage{xcolor}
\usepackage{microtype}
\usepackage{booktabs}
\usetikzlibrary{calc}

% --- Definizione Colori (codice) ---
\definecolor{backcolour}{RGB}{250, 250, 250}
\definecolor{framegray}{RGB}{220, 220, 220}
\definecolor{codeblue}{RGB}{0, 0, 139}
\definecolor{codegreen}{RGB}{0, 100, 100}
\definecolor{codegray}{RGB}{128, 128, 128}
\definecolor{numgray}{RGB}{180, 180, 180}

% --- Stile Codice ---
\lstdefinestyle{exactstyle}{
    language=C,
    backgroundcolor=\color{backcolour},
    commentstyle=\upshape\color{codegray},
    keywordstyle=\bfseries\color{codeblue},
    stringstyle=\color{codegreen},
    basicstyle=\ttfamily\small\upshape,
    breakatwhitespace=false,
    breaklines=true,
    keepspaces=true,
    numbers=left,
    numbersep=15pt,
    numberstyle=\tiny\color{numgray},
    showspaces=false,
    showstringspaces=false,
    showtabs=false,
    tabsize=4,
    frame=none,
    xleftmargin=25pt,
}
\lstset{style=exactstyle}

% --- Box Strategia ---
\newtcolorbox{strategiabox}{
  colback=white!3,
  colframe=black,
  boxrule=0.8pt,
  title=\textbf{Metodologia risolutiva proposta:},
  fonttitle=\bfseries,
  colbacktitle=white!15,
  coltitle=black
}
\newcommand{\strategia}[1]{%
  \begin{strategiabox}
    #1
  \end{strategiabox}
}

% --- Theorem-like ambienti semplici ---
\newtheorem{esempio}{Esempio}[chapter]
\newtheorem{osservazione}{Osservazione}[chapter]

% --- Colori per i box dei teoremi ---
\definecolor{defscol}{HTML}{ecd8d7}
\definecolor{asumscol}{HTML}{ecd8d7}
\definecolor{rmkscol}{HTML}{313160}
\definecolor{exmscol}{HTML}{e04b52}
\definecolor{lemscol}{HTML}{2c3943}
\definecolor{thmscol}{HTML}{595765}
\definecolor{prpscol}{HTML}{9c98b1}
\definecolor{corscol}{HTML}{dfd9fd}
\definecolor{clmscol}{HTML}{165c58}
\definecolor{facscol}{HTML}{28a8a1}

% ============================
% Definizione
% ============================
\newtcbtheorem[number within=section]{mydefinition}{Definizione}{
    enhanced, frame hidden, titlerule=0mm, toptitle=1mm, bottomtitle=1mm,
    fonttitle=\bfseries\large, coltitle=black,
    colbacktitle=defscol!40!white, colback=defscol!20!white,
}{defn}
\NewDocumentCommand{\defn}{m+m}{\begin{mydefinition}{#1}{}#2\end{mydefinition}}
\NewDocumentCommand{\defnr}{mm+m}{\begin{mydefinition}{#1}{#2}#3\end{mydefinition}}

% ============================
% Proprietà
% ============================
\newtcbtheorem[number within=section]{mypropi}{Proprietà}{
    enhanced, frame hidden, titlerule=0mm, toptitle=1mm, bottomtitle=1mm,
    fonttitle=\bfseries\large, coltitle=black,
    colbacktitle=blue!40!white, colback=blue!20!white,
}{propi}
\NewDocumentCommand{\propi}{m+m}{
    \begin{mypropi}{#1}{}\begin{itemize}#2\end{itemize}\end{mypropi}
}

% ============================
% Assunzione
% ============================
\newtcbtheorem[use counter from=mydefinition]{myassumption}{Assunzione}{
    enhanced, frame hidden, titlerule=0mm, toptitle=1mm, bottomtitle=1mm,
    fonttitle=\bfseries\large, coltitle=black,
    colbacktitle=asumscol!40!white, colback=asumscol!20!white,
}{asum}
\NewDocumentCommand{\asum}{m+m}{\begin{myassumption}{#1}{}#2\end{myassumption}}
\NewDocumentCommand{\asumr}{mm+m}{\begin{myassumption}{#1}{#2}#3\end{myassumption}}

% ============================
% Teorema
% ============================
\newtcbtheorem[number within=section]{mytheorem}{Teorema}{
    enhanced, frame hidden, titlerule=0mm, toptitle=1mm, bottomtitle=1mm,
    fonttitle=\bfseries\large, coltitle=black,
    colbacktitle=thmscol!40!white, colback=thmscol!20!white,
}{thm}
\NewDocumentCommand{\thm}{m+m}{\begin{mytheorem}{#1}{}#2\end{mytheorem}}
\NewDocumentCommand{\thmr}{mm+m}{\begin{mytheorem}{#1}{#2}#3\end{mytheorem}}

\newenvironment{thmpf}{
    {\noindent{\it \textbf{Dimostrazione.}}}
    \tcolorbox[blanker,breakable,left=5mm,parbox=false,
    before upper={\parindent15pt}, after skip=10pt,
    borderline west={1mm}{0pt}{thmscol!40!white}]
}{
    \textcolor{thmscol!40!white}{\hbox{}\nobreak\hfill$\blacksquare$}
    \endtcolorbox
}
\NewDocumentCommand{\thmp}{m+m+m}{
    \begin{mytheorem}{#1}{}#2\end{mytheorem}
    \begin{thmpf}#3\end{thmpf}
}

% ============================
% Lemma
% ============================
\newtcbtheorem[use counter from=mydefinition]{mylemma}{Lemma}{
    enhanced, frame hidden, titlerule=0mm, toptitle=1mm, bottomtitle=1mm,
    fonttitle=\bfseries\large, coltitle=black,
    colbacktitle=lemscol!40!white, colback=lemscol!20!white,
}{lem}
\NewDocumentCommand{\lem}{m+m}{\begin{mylemma}{#1}{}#2\end{mylemma}}

\newenvironment{lempf}{
    {\noindent{\it \textbf{Dimostrazione del Lemma.}}}
    \tcolorbox[blanker,breakable,left=5mm,parbox=false,
    before upper={\parindent15pt}, after skip=10pt,
    borderline west={1mm}{0pt}{lemscol!40!white}]
}{
    \textcolor{lemscol!40!white}{\hbox{}\nobreak\hfill$\blacksquare$}
    \endtcolorbox
}
\NewDocumentCommand{\lemp}{m+m+m}{
    \begin{mylemma}{#1}{}#2\end{mylemma}
    \begin{lempf}#3\end{lempf}
}

% ============================
% Corollario
% ============================
\newtcbtheorem[use counter from=mydefinition]{mycorollary}{Corollario}{
    enhanced, frame hidden, titlerule=0mm, toptitle=1mm, bottomtitle=1mm,
    fonttitle=\bfseries\large, coltitle=black,
    colbacktitle=corscol!40!white, colback=corscol!20!white,
}{cor}
\NewDocumentCommand{\cor}{+m}{\begin{mycorollary}{}{}#1\end{mycorollary}}

\newenvironment{corpf}{
    {\noindent{\it \textbf{Dimostrazione del Corollario.}}}
    \tcolorbox[blanker,breakable,left=5mm,parbox=false,
    before upper={\parindent15pt}, after skip=10pt,
    borderline west={1mm}{0pt}{corscol!40!white}]
}{
    \textcolor{corscol!40!white}{\hbox{}\nobreak\hfill$\blacksquare$}
    \endtcolorbox
}
\NewDocumentCommand{\corp}{m+m+m}{
    \begin{mycorollary}{}{}#1\end{mycorollary}
    \begin{corpf}#2\end{corpf}
}

% ============================
% Proposizione
% ============================
\newtcbtheorem[use counter from=mydefinition]{myproposition}{Proposizione}{
    enhanced, frame hidden, titlerule=0mm, toptitle=1mm, bottomtitle=1mm,
    fonttitle=\bfseries\large, coltitle=black,
    colbacktitle=prpscol!30!white, colback=prpscol!20!white,
}{prop}
\NewDocumentCommand{\prop}{+m}{\begin{myproposition}{}{}#1\end{myproposition}}

\newenvironment{proppf}{
    {\noindent{\it \textbf{Dimostrazione.}}}
    \tcolorbox[blanker,breakable,left=5mm,parbox=false,
    before upper={\parindent15pt}, after skip=10pt,
    borderline west={1mm}{0pt}{prpscol!40!white}]
}{
    \textcolor{prpscol!40!white}{\hbox{}\nobreak\hfill$\blacksquare$}
    \endtcolorbox
}
\NewDocumentCommand{\propp}{+m+m}{
    \begin{myproposition}{}{}#1\end{myproposition}
    \begin{proppf}#2\end{proppf}
}

% ============================
% Claim
% ============================
\newtcbtheorem[use counter from=mydefinition]{myclaim}{Claim}{
    enhanced, frame hidden, titlerule=0mm, toptitle=1mm, bottomtitle=1mm,
    fonttitle=\bfseries\large, coltitle=black,
    colbacktitle=clmscol!40!white, colback=clmscol!20!white,
}{clm}
\NewDocumentCommand{\clm}{m+m}{\begin{myclaim*}{#1}{}#2\end{myclaim*}}

\newenvironment{clmpf}{
    {\noindent{\it \textbf{Dimostrazione del Claim.}}}
    \tcolorbox[blanker,breakable,left=5mm,parbox=false,
    before upper={\parindent15pt}, after skip=10pt,
    borderline west={1mm}{0pt}{clmscol!40!white}]
}{
    \textcolor{clmscol!40!white}{\hbox{}\nobreak\hfill$\blacksquare$}
    \endtcolorbox
}
\NewDocumentCommand{\clmp}{m+m+m}{
    \begin{myclaim*}{#1}{}#2\end{myclaim*}
    \begin{clmpf}#3\end{clmpf}
}

% ============================
% Fatto
% ============================
\newtcbtheorem[use counter from=mydefinition]{myfact}{Fatto}{
    enhanced, frame hidden, titlerule=0mm, toptitle=1mm, bottomtitle=1mm,
    fonttitle=\bfseries\large, coltitle=black,
    colbacktitle=facscol!40!white, colback=facscol!20!white,
}{fact}
\NewDocumentCommand{\fact}{+m}{\begin{myfact}{}{}#1\end{myfact}}

% ============================
% Dimostrazione semplice
% ============================
\NewDocumentCommand{\pf}{+m}{
    \begin{proof}[\noindent\textbf{Dimostrazione.}]#1\end{proof}
}

% ============================
% Esempio (box)
% ============================
\newenvironment{myexample}{
    \tcolorbox[blanker,breakable,left=5mm,parbox=false,
    before upper={\parindent15pt}, after skip=10pt,
    borderline west={1mm}{0pt}{clmscol!40!white}]
}{
    \textcolor{clmscol!40!white}{\hbox{}\nobreak\hfill$\blacksquare$}
    \endtcolorbox
}
\NewDocumentCommand{\exm}{m+m}{
    \begin{myexample}
    {\noindent{\it \textbf{Dimostrazione: #1}}}\\
        #2
    \end{myexample}
}

% ============================
% Remark / Spiegazione
% ============================
\NewDocumentCommand{\rmk}{+m}{
    {\it \color{rmkscol!40!white}#1}
}
\newenvironment{remark}{
    \par\vspace{5pt}
    {\par\noindent{\textbf{Spiegazione.}}}
    \tcolorbox[blanker,breakable,left=5mm,
    before skip=10pt,after skip=10pt,
    borderline west={1mm}{0pt}{rmkscol!20!white}]
}{
    \endtcolorbox
    \vspace{5pt}
}
\NewDocumentCommand{\rmkb}{+m}{\begin{remark}#1\end{remark}}

% --- Comandi extra ---
\makeatletter
\newcommand{\impliesleft}{\mathrel{\mathpalette\implies@left\relax}}
\newcommand{\implies@left}[2]{\reflectbox{$\m@th#1\implies$}}
\makeatother

% --- Titlepage personalizzata (sovrascrive \maketitle di Quarto) ---
\renewcommand{\maketitle}{%
  \begin{titlepage}
    \vspace*{\fill}
    \centering
    {\LARGE \textbf{Basi di Dati — Modulo 2}\par}
    \vspace{1cm}
    {\Large Anno Accademico 2025-2026\par}
    \vspace{1cm}
    {\normalsize Dispensa a cura di\\[2mm] Ibrahima Tely Barry \par}
    \vspace*{\fill}
  \end{titlepage}
}
